\section{The projection in the prescribed ultrapower}\label{r:sec:completion}

We now pass from the finite-set result to one projection, keeping the
designated trace family and the fixed ultrafilter throughout.

\begin{lemma}[Detection by limit traces]\label{r:completion:detection}
For a bounded sequence \(d=(d_j)\) in \(B\), with class
\([d]\in B^\omega\), one has
\begin{align}
 \sup_{\lambda\in\Lambda_\omega}\lambda([d]^*[d])
 &=\lim_{j\to\omega}\sup_{\tau\in\T(A)}
             \bar\tau(d_j^*d_j),\label{r:eq:limit-two}\\
 \sup_{\lambda\in\Lambda_\omega}|\lambda([d])|
 &=\lim_{j\to\omega}\sup_{\tau\in\T(A)}
             |\bar\tau(d_j)|.\label{r:eq:limit-linear}
\end{align}
In particular, if \(\normtwo z\lambda=0\) for every
\(\lambda\in\Lambda_\omega\), then \(z=0\) in \(B^\omega\).
\end{lemma}

\begin{proof}
The upper bounds follow from the definition of a limit trace.
For the reverse bound in \eqref{r:eq:limit-two}, choose \(\tau_j\) whose
value on \(d_j^*d_j\) is within \(1/j\) of the supremum.
The associated limit trace gives equality.
For \eqref{r:eq:limit-linear}, choose approximate maximizers of
\(|\bar\tau(d_j)|\). Absolute value commutes with the ultralimit of a
bounded complex sequence, so this choice also gives equality.
The final assertion is \eqref{r:eq:limit-two} and the definition of
\(I_\omega\).
\end{proof}

\begin{proof}[Proof of Theorem~\ref{r:thm:main}]
Set \(D=B^\omega\). It is a unital real-rank-zero \(C^*\)-algebra by
hypothesis and has the tracial states in \(\Lambda_\omega\).
The inclusion \(A\subseteq D\) is unital and faithful.
Lemma~\ref{r:proj:small} supplies full projections of uniformly
arbitrarily small trace over all \(\T(D)\).
Thus Theorem~\ref{r:thm:finite} applies.

Choose increasing finite subsets \(F_n\subseteq\U(A)\) whose union is
norm dense in \(\U(A)\), including \(1\) in each \(F_n\).
Theorem~\ref{r:thm:finite} gives projections \(P_n\in D\) with
\begin{equation}\label{r:eq:finite-stage}
 \normtwo{[P_n,x]}{\tau}<1/n,\qquad
 \left|\tau(P_nx)-\tfrac12\tau(x)\right|<1/n
 \quad(x\in F_n,\ \tau\in\T(D)).
\end{equation}
Lift \(P_n\) to positive contractions \((c_{n,j})_j\) in \(B\).
Such lifts exist by first taking a self-adjoint lift and then applying
the continuous function which clips the real line to \([0,1]\).
Since \(P_n\) is a projection,
\[
 \lim_{j\to\omega}\normtwo{c_{n,j}^2-c_{n,j}}{\T(A)}=0.
\]

For each fixed \(n\), Lemma~\ref{r:completion:detection}, applied to the
commutators and the linear errors in \eqref{r:eq:finite-stage}, gives
an \(\omega\)-large set of coordinates on which
\begin{align*}
 \normtwo{c_{n,j}^2-c_{n,j}}{\T(A)}&<1/n,\\
 \normtwo{[c_{n,j},x]}{\T(A)}&<2/n &&(x\in F_n),\\
 \sup_{\tau\in\T(A)}
 \left|\bar\tau(c_{n,j}x)-\tfrac12\bar\tau(x)\right|
 &<2/n &&(x\in F_n).
\end{align*}
There are finitely many conditions, so their good coordinate sets
have an \(\omega\)-large intersection.
Choose one coordinate \(j(n)\) in that intersection and put
\(b_n=c_{n,j(n)}\).

The sequence \((b_n)\) consists of positive contractions.
Its projection defect, and every eventually included test error,
converge to zero in the ordinary sense. Hence
\(p=[(b_n)]\in B^\omega\) is a projection, commutes with the tested
unitaries, and satisfies
\(\lambda(px)=\lambda(x)/2\) on those unitaries for every
\(\lambda\in\Lambda_\omega\).
Norm density and linear spanning by unitaries extend both conclusions
to every \(x\in A\).
This construction has used the original \(\omega\); no further
ultrapower is taken.

Finally, \(A\) is dense in \(B\) for the designated uniform tracial
2-norm. Indeed, if \(x=[(a_j)]\in C/J=B\), then
\[
 \normtwo{x-a_m}{\T(A)}^2
 =\sup_{\tau\in\T(A)}
       \lim_j\tau((a_j-a_m)^*(a_j-a_m))\longrightarrow0
\]
by uniform 2-Cauchyness.
Given \(x\in B\), choose \(a\in A\) with
\(\normtwo{x-a}{\T(A)}<\delta\).
For every \(\lambda\in\Lambda_\omega\), multiplication by the
contraction \(p\) and tracial Cauchy--Schwarz give
\[
 \normtwo{[p,x]}{\lambda}\leq2\delta,\qquad
 \left|\lambda(px)-\tfrac12\lambda(x)\right|\leq\tfrac32\delta.
\]
Let \(\delta\) tend to zero. The first assertion, together with
Lemma~\ref{r:completion:detection}, gives \([p,x]=0\) in \(D\) itself;
the second is the required trace identity.
Thus the same \(p\) works for the entire constant copy of \(B\), even when
\(B\) is not norm separable.
\end{proof}

We finish by identifying all abstract traces of the completion, as stated
in Corollary~\ref{r:cor:completion-traces}.

\begin{proof}[Proof of Corollary~\ref{r:cor:completion-traces}]
Put $X_B=\{\bar\tau:\tau\in T(A)\}$. By
\cite[Proposition~3.23(ii)--(iv)]{CCEGSTW}, $X_B$ is a compact face of
$T(B)$ and the designated uniform $2$-norm is complete on the unit ball
of $B$. In the terminology of that paper, $(B,X_B)$ is a factorial
tracially complete pair.
Theorem~\ref{r:thm:main} and
\cite[Proposition~5.20]{CCEGSTW} give property $\Gamma$ for this pair,
and hence complemented partitions of unity by
\cite[Theorem~1.4]{CCEGSTW}. Evington's trace theorem
\cite[Theorem~B]{EvingtonTraces25} then gives $T(B)=X_B$ provided the
pair is of type $\mathrm{II}_1$: each designated GNS von Neumann algebra
must have no nonzero type I central summand.

For every $\tau\in T(A)$, the finite von Neumann algebra
$\pi_\tau(A)''$ has no nonzero type I central summand. Otherwise it
would have a nonzero homogeneous summand $M_n(Z)$ with $Z$ abelian;
composing the representation of $A$ into this summand with a character
of $Z$ would give a unital finite-dimensional representation of $A$.
Simplicity would make that representation injective, contrary to the
infinite dimensionality of $A$. Thus every such GNS algebra is of
type $\mathrm{II}_1$. Their preservation under completion
\cite[Proposition~3.23(vi)]{CCEGSTW} makes $(B,X_B)$ a type
$\mathrm{II}_1$ pair, so the cited trace theorem applies.
Restriction is therefore a continuous affine bijection between compact
Hausdorff trace spaces, and hence a homeomorphism. Finally, for
$\sigma=\bar\tau$, tracial Cauchy--Schwarz gives
$|\sigma(b)|\le\bar\tau(b^*b)^{1/2}\le\|b\|_{2,T(A)}$.
\end{proof}
